% GENERATED FILE. Edit canonical lessons, then run npm run generate:books.
% canonical-source-hash: fnv1a64:6d5518a84c21d1c7
% canonical-lessons: LA-C522-palleo, LA-C522-pallidus, LA-C522-rubeo, LA-C522-macer, LA-C522-macies

\chapter{Pale, red and thin}
\label{ch:la-pale-red-and-thin}
\hlchaptermodality{522}

\begin{chapteropening}
\textbf{By the end of this chapter:} \emph{I can describe how someone looks when unwell.}

Complete the last lesson of chapter 522: I can describe how someone looks when unwell.
\end{chapteropening}

\section[palleō, pallēre]{palleō, pallēre — to be pale}

\label{lesson:LA-C522-palleo}



\begin{quote}
Before the new one: say the Latin for weak, then the Latin for weakness.
\end{quote}

\subsection*{You'll want to know: palleō, pallēre}
\textbf{palleō, pallēre} — \textquotedblleft{}to be pale\textquotedblright{}.

\begin{quote}
Ancilla timōre pallet.
\end{quote}

\begin{quote}
\emph{The maid is pale with fear.}
\end{quote}

\emph{palleō, pallēre, palluī}. English \emph{pallor} comes from its noun.

\begin{cousinweb}
\textbf{palleō} is inherited from a Proto-Indo-European root meaning \textquotedblleft{}grey, pale\textquotedblright{} (EtymDB; Lewis \& Short). English \emph{pallor} comes from the noun built on it.
\end{cousinweb}

\begin{grammarlens}[title={What you've built}]
One more word for telling, explaining and describing at B1.
\end{grammarlens}

\subsection*{Your turn}
\begin{itemize}
\raggedright
  \item \emph{Say it:} \emph{palleō, pallēre}
  \item \emph{Say it:} \emph{palleō, pallēre}, once more
  \item \emph{Say it:} a sentence of your own with \emph{palleō, pallēre}
  \item \emph{Recall:} say the Latin for weak, then the Latin for weakness, then say \emph{palleō, pallēre} again
\end{itemize}

\begin{tcolorbox}[breakable,colback=teal!4,colframe=teal!35!black,title={Before you move on}]
What does palleō, pallēre mean? (\textquotedblleft{}To be pale\textquotedblright{}.) Say it once more.
\end{tcolorbox}
\section[pallidus, pallida, pallidum]{pallidus, pallida, pallidum — pale}

\label{lesson:LA-C522-pallidus}



\begin{quote}
Before the new one: say the Latin for weakness, then the Latin for to be pale.
\end{quote}

\subsection*{You'll want to know: pallidus, pallida, pallidum}
\textbf{pallidus, pallida, pallidum} — \textquotedblleft{}pale\textquotedblright{}.

\begin{quote}
Post morbum puella pallida est.
\end{quote}

\begin{quote}
\emph{After her illness the girl is pale.}
\end{quote}

From \emph{palleō}. English \emph{pallid}, and \emph{pale} through French.

\begin{cousinweb}
\textbf{pallidus} is built on \emph{palleō} \textquotedblleft{}to be pale\textquotedblright{} with \emph{-idus} (Lewis \& Short; EtymDB). English \emph{pallid} comes from it, and \emph{pale} came the same way through French.
\end{cousinweb}

\begin{grammarlens}[title={What you've built}]
One more word for telling, explaining and describing at B1.
\end{grammarlens}

\subsection*{Your turn}
\begin{itemize}
\raggedright
  \item \emph{Say it:} \emph{pallidus, pallida, pallidum}
  \item \emph{Say it:} \emph{pallidus, pallida, pallidum}, once more
  \item \emph{Say it:} a sentence of your own with \emph{pallidus, pallida, pallidum}
  \item \emph{Recall:} say the Latin for weakness, then the Latin for to be pale, then say \emph{pallidus, pallida, pallidum} again
\end{itemize}

\begin{tcolorbox}[breakable,colback=teal!4,colframe=teal!35!black,title={Before you move on}]
What does pallidus, pallida, pallidum mean? (\textquotedblleft{}Pale\textquotedblright{}.) Say it once more.
\end{tcolorbox}
\section[rubeō, rubēre]{rubeō, rubēre — to be red}

\label{lesson:LA-C522-rubeo}



\begin{quote}
Before the new one: say the Latin for to be pale, then the Latin for pale.
\end{quote}

\subsection*{You'll want to know: rubeō, rubēre}
\textbf{rubeō, rubēre} — \textquotedblleft{}to be red\textquotedblright{}.

\begin{quote}
Genae puerī rubent.
\end{quote}

\begin{quote}
\emph{The boy's cheeks are red.}
\end{quote}

\emph{rubeō, rubēre, rubuī}, from \emph{ruber} — red. \emph{Ērubēscō}, already known, is to turn red.

\begin{cousinweb}
\textbf{rubeō} belongs with \emph{ruber} \textquotedblleft{}red\textquotedblright{}, inherited from a Proto-Indo-European root meaning \textquotedblleft{}red\textquotedblright{} (EtymDB). English \emph{ruby} and \emph{rubicund} come from the same family, and English \emph{red} is a cousin.
\end{cousinweb}

\begin{grammarlens}[title={What you've built}]
One more word for telling, explaining and describing at B1.
\end{grammarlens}

\subsection*{Your turn}
\begin{itemize}
\raggedright
  \item \emph{Say it:} \emph{rubeō, rubēre}
  \item \emph{Say it:} \emph{rubeō, rubēre}, once more
  \item \emph{Say it:} a sentence of your own with \emph{rubeō, rubēre}
  \item \emph{Recall:} say the Latin for to be pale, then the Latin for pale, then say \emph{rubeō, rubēre} again
\end{itemize}

\begin{tcolorbox}[breakable,colback=teal!4,colframe=teal!35!black,title={Before you move on}]
What does rubeō, rubēre mean? (\textquotedblleft{}To be red\textquotedblright{}.) Say it once more.
\end{tcolorbox}
\section[macer, macra, macrum]{macer, macra, macrum — thin, lean}

\label{lesson:LA-C522-macer}



\begin{quote}
Before the new one: say the Latin for pale, then the Latin for to be red.
\end{quote}

\subsection*{You'll want to know: macer, macra, macrum}
\textbf{macer, macra, macrum} — \textquotedblleft{}thin, lean\textquotedblright{}.

\begin{quote}
Canis macer cibum quaerit.
\end{quote}

\begin{quote}
\emph{The thin dog looks for food.}
\end{quote}

Like \emph{ruber}: \emph{macer}, \emph{macra}, \emph{macrum}. English \emph{meagre} comes from it through French.

\begin{cousinweb}
\textbf{macer} is inherited from a Proto-Indo-European word meaning \textquotedblleft{}thin\textquotedblright{} (EtymDB), and German \emph{mager} \textquotedblleft{}lean\textquotedblright{} is a cousin (Lewis \& Short). English \emph{meagre} comes from it through French.
\end{cousinweb}

\begin{grammarlens}[title={What you've built}]
One more word for telling, explaining and describing at B1.
\end{grammarlens}

\subsection*{Your turn}
\begin{itemize}
\raggedright
  \item \emph{Say it:} \emph{macer, macra, macrum}
  \item \emph{Say it:} \emph{macer, macra, macrum}, once more
  \item \emph{Say it:} a sentence of your own with \emph{macer, macra, macrum}
  \item \emph{Recall:} say the Latin for pale, then the Latin for to be red, then say \emph{macer, macra, macrum} again
\end{itemize}

\begin{tcolorbox}[breakable,colback=teal!4,colframe=teal!35!black,title={Before you move on}]
What does macer, macra, macrum mean? (\textquotedblleft{}Thin, lean\textquotedblright{}.) Say it once more.
\end{tcolorbox}
\section[maciēs, maciēī]{maciēs, maciēī — leanness, thinness}

\label{lesson:LA-C522-macies}



\begin{quote}
Before the new one: say the Latin for to be red, then the Latin for thin.
\end{quote}

\subsection*{You'll want to know: maciēs, maciēī}
\textbf{maciēs, maciēī} — \textquotedblleft{}leanness, thinness\textquotedblright{}.

\begin{quote}
Maciēs eius morbum ostendit.
\end{quote}

\begin{quote}
\emph{His thinness shows his illness.}
\end{quote}

Feminine, fifth declension. English \emph{emaciated} comes from the related verb \emph{ēmaciō}.

\begin{cousinweb}
\textbf{maciēs} belongs with \emph{macer} \textquotedblleft{}thin\textquotedblright{} (EtymDB; Lewis \& Short). English \emph{emaciated} comes from the related verb \emph{ēmaciō}.
\end{cousinweb}

\begin{grammarlens}[title={What you've built}]
One more word for telling, explaining and describing at B1.
\end{grammarlens}

\subsection*{Your turn}
\begin{itemize}
\raggedright
  \item \emph{Say it:} \emph{maciēs, maciēī}
  \item \emph{Say it:} \emph{maciēs, maciēī}, once more
  \item \emph{Say it:} a sentence of your own with \emph{maciēs, maciēī}
  \item \emph{Recall:} say the Latin for to be red, then the Latin for thin, then say \emph{maciēs, maciēī} again
\end{itemize}

\begin{tcolorbox}[breakable,colback=teal!4,colframe=teal!35!black,title={Before you move on}]
What does maciēs, maciēī mean? (\textquotedblleft{}Leanness, thinness\textquotedblright{}.) Say it once more.
\end{tcolorbox}
